% Options for packages loaded elsewhere
\PassOptionsToPackage{unicode}{hyperref}
\PassOptionsToPackage{hyphens}{url}
\PassOptionsToPackage{dvipsnames,svgnames,x11names}{xcolor}
%
\documentclass[
  11pt,
  a4paper,
]{ltjsarticle}
\usepackage{amsmath,amssymb}
\usepackage{iftex}
\ifPDFTeX
  \usepackage[T1]{fontenc}
  \usepackage[utf8]{inputenc}
  \usepackage{textcomp} % provide euro and other symbols
\else % if luatex or xetex
  \usepackage{unicode-math} % this also loads fontspec
  \defaultfontfeatures{Scale=MatchLowercase}
  \defaultfontfeatures[\rmfamily]{Ligatures=TeX,Scale=1}
\fi
\usepackage{lmodern}
\ifPDFTeX\else
  % xetex/luatex font selection
\fi
% Use upquote if available, for straight quotes in verbatim environments
\IfFileExists{upquote.sty}{\usepackage{upquote}}{}
\IfFileExists{microtype.sty}{% use microtype if available
  \usepackage[]{microtype}
  \UseMicrotypeSet[protrusion]{basicmath} % disable protrusion for tt fonts
}{}
\makeatletter
\@ifundefined{KOMAClassName}{% if non-KOMA class
  \IfFileExists{parskip.sty}{%
    \usepackage{parskip}
  }{% else
    \setlength{\parindent}{0pt}
    \setlength{\parskip}{6pt plus 2pt minus 1pt}}
}{% if KOMA class
  \KOMAoptions{parskip=half}}
\makeatother
\usepackage{xcolor}
\usepackage[margin=24mm]{geometry}
\setlength{\emergencystretch}{3em} % prevent overfull lines
\providecommand{\tightlist}{%
  \setlength{\itemsep}{0pt}\setlength{\parskip}{0pt}}
\setcounter{secnumdepth}{5}
\ifLuaTeX
  \usepackage{selnolig}  % disable illegal ligatures
\fi
\IfFileExists{bookmark.sty}{\usepackage{bookmark}}{\usepackage{hyperref}}
\IfFileExists{xurl.sty}{\usepackage{xurl}}{} % add URL line breaks if available
\urlstyle{same}
\hypersetup{
  colorlinks=true,
  linkcolor={blue},
  filecolor={Maroon},
  citecolor={Blue},
  urlcolor={blue},
  pdfcreator={LaTeX via pandoc}}

\author{}
\date{}

\begin{document}

\hypertarget{qr-ux5206ux89e3}{%
\section{QR 分解}\label{qr-ux5206ux89e3}}

行列 \texttt{A∈R\^{}\{m×n\}} を

\[
A=QR
\]

と分解します。\texttt{Q} の列は正規直交し、\texttt{R} は上三角行列です。

\hypertarget{ux5e7eux4f55ux5b66ux7684ux610fux5473}{%
\subsection{幾何学的意味}\label{ux5e7eux4f55ux5b66ux7684ux610fux5473}}

\texttt{A} の列が張る部分空間は \texttt{Q}
の列が張る部分空間と同じです。しかし \texttt{Q}
は正規直交基底なので、座標計算が非常に簡単になります。

\texttt{A} の複雑さを

\begin{itemize}
\tightlist
\item
  \texttt{Q}: 空間の向き
\item
  \texttt{R}: その基底での係数
\end{itemize}

へ分離していると理解できます。

\hypertarget{ux6700ux5c0fux4e8cux4e57}{%
\subsection{最小二乗}\label{ux6700ux5c0fux4e8cux4e57}}

\[
\min_x\|Ax-b\|_2
\]

に \texttt{A=QR} を代入すると

\[
\|QRx-b\|_2
\]

となります。\texttt{Q} の直交性を利用すると、正規方程式で
\texttt{A\^{}TA} を作るより数値的に安定に解けます。

\hypertarget{gramschmidt-ux3068ux306eux95a2ux4fc2}{%
\subsection{Gram--Schmidt
との関係}\label{gramschmidt-ux3068ux306eux95a2ux4fc2}}

Gram--Schmidt は QR を構成する方法の一つです。ただし数値計算では
Householder 変換など、より安定な方法が一般に使われます。

\end{document}
